% titel: Blatt 8, A3: Flächeninhalte und Stammfunktionen
% skript: Kap. 5 (Stammfunktionen, Hauptsatz)

% Teil I: Werkzeuge, kein Lösungsweg.

Stammfunktionen zu verschiedenen Funktionen angeben, Flächen zwischen zwei
Kurven bzw. zwischen Kurve und $x$-Achse berechnen und einen Grenzwert für
unendliche Integrationsgrenze auswerten.

\begin{werkzeugkasten}[Grund-Stammfunktionen]
\begin{itemize}
  \item $\displaystyle\int\frac1x\,dx=\ln|x|+\mathbb C$ ($x\ne0$).
  \item Mit der Kettenregel-Rückwärts (Substitution): Aus
        $\int u'(x)u(x)\,dx$ Typen wie
        $2\sin x\cos x=\sin(2x)$ erkennen und integrieren.
  \item $\displaystyle\frac{2x}{\cos^2(x^2)}$ ist die Ableitung von
        $\tan(x^2)$ (Kettenregel in der Form
        $[\tan g]'=\frac{g'}{\cos^2g}$ mit $g=x^2$, $g'=2x$).
\end{itemize}
\textit{Skript:} Kap. 5 (Stammfunktionen).
\end{werkzeugkasten}

\begin{werkzeugkasten}[Fläche zwischen zwei Kurven]
Die von $f$ und $g$ auf $[a,b]$ eingeschlossene Fläche ist
\begin{equation*}
  \int_a^b |f(x)-g(x)|\,dx .
\end{equation*}
Hat $f-g$ ein festes Vorzeichen auf Teilintervallen, rechnet man die Abschnitte
getrennt. Beispielhaft für $f(x)=\sin x$ und $g(x)=-x(x-\pi)(x-2\pi)$ auf
$[0,2\pi]$ erhält man über die Stammfunktionen von $\sin$ und des
kubischen Polynoms einen exakten Ausdruck.

\textit{Skript:} Kap. 5 (Hauptsatz, Flächen).
\end{werkzeugkasten}

\begin{achtung}[Darauf achten]
\begin{itemize}
  \item Betrag nicht vergessen: $\int(f-g)$ ist nicht die Fläche, falls
        $f<g$ auf einem Teilbereich.
  \item Bei $1/x$ unbedingt $\ln|x|$ statt $\ln x$ schreiben (Integrand
        auch für $x<0$ erklärt).
  \item Skizzieren zuerst: die Vorzeichen der Differenz $f-g$ bestimmen
        die Zerlegung der Fläche.
  \item Für die Grenzwert-Aufgabe c) die Stammfunktion exakt hinschreiben,
        die Grenzen einsetzen und den Grenzwert $b\to\infty$ bilden.
\end{itemize}
\end{achtung}
